\chapter{System Analysis and Design}
\label{ch:analysis}

\section{Introduction}

Getting from stakeholder pain points to a working system design is rarely a linear process. This chapter records what that process produced for the IPDC platform: the requirements that emerged from stakeholder analysis, the framework for adapting Chinese smart park patterns to the Ethiopian context, and the architecture that resulted. Not every decision was obvious at the time --- but all of them trace back to a documented constraint or a real operational need.

\section{Requirements Analysis}

\subsection{Stakeholder Identification and Analysis}

Organizational analysis of IPDC's structure, combined with insights from the stakeholder questionnaires, identified four distinct stakeholder groups, each with its own role, priorities, and interaction patterns.

\textbf{Tenant Enterprises} are the primary consumers of services. Their representatives --- typically operations managers and administrative staff --- submit service requests, track progress, handle fee payments, and raise complaints. Speed, transparency, and the ability to keep working during connectivity outages are their chief concerns.

\textbf{IPDC Administrative Staff} serve as the bridge between tenants and operational resources. Service coordinators, permit officers, and finance officers process requests, publish announcements, manage tenant accounts, and resolve complaints. Workflow efficiency and dependable record-keeping matter most to this group.

\textbf{Operations Personnel} are the ones doing the physical work once a request is approved. A maintenance coordinator might receive a task assignment, find they cannot access the system from a remote corner of the park, and need to pick it back up an hour later. That specific scenario --- task received, connectivity lost, work continues --- is exactly what the mobile offline task queue was built for.

\textbf{IPDC Management} needs a high-level view of operational performance without involvement in day-to-day platform use. Directors and planning officers rely on aggregated data about service volumes, resolution times, and resource utilization to inform decisions.

\subsection{Functional Requirements}

Table~\ref{tab:func-req} summarizes the twelve functional requirement groups that emerged from stakeholder analysis and the mapping of Chinese smart park features.

\begin{table}[H]
\centering
\caption{Functional Requirements Summary}
\label{tab:func-req}
\small
\begin{tabular}{lp{3.5cm}p{7.5cm}}
\toprule
\textbf{ID} & \textbf{Requirement} & \textbf{Description} \\
\midrule
FR1 & User Authentication & Email-based registration and login with role assignment (tenant, admin, operations); offline authentication fallback using cached credentials \\
FR2 & Service Request Management & Submit, track, approve, and complete service requests across multiple categories with file attachments and priority levels \\
FR3 & Digital Service Token System & Internal credit system for fee tracking with tiered accounts (basic, silver, gold, platinum), transaction history, and automated deductions \\
FR4 & Announcement Management & Create, distribute, and display system announcements targeted to specific audiences with priority and expiration settings \\
FR5 & Complaint Management & File complaints against completed services with evidence uploads, track resolution through defined workflows \\
FR6 & Asset \& Maintenance Management & Register, categorize, and track park assets; schedule and record maintenance activities \\
FR7 & AI-Powered Services & Automated service request classification, predictive equipment maintenance, and intelligent park recommendation for prospective tenants \\
FR8 & Document Management & Upload, store, and retrieve documents associated with service requests and OSS applications \\
FR9 & OSS Services & Digital interfaces for investment permits, business licenses, customs services, and banking coordination \\
FR10 & Interactive Park Map & Geographic visualization of Ethiopian industrial parks with detail overlays \\
FR11 & Billing \& Invoicing & Generate invoices, track payment transactions, and export PDF reports \\
FR12 & Notifications & In-platform and email notifications for status changes and announcements \\
\bottomrule
\end{tabular}
\end{table}

\subsection{Non-Functional Requirements}

Table~\ref{tab:nonfunc-req} sets out the non-functional requirements that constrain and shape the system architecture.

\begin{table}[H]
\centering
\caption{Non-Functional Requirements}
\label{tab:nonfunc-req}
\small
\begin{tabular}{lp{3cm}p{8cm}}
\toprule
\textbf{ID} & \textbf{Requirement} & \textbf{Specification} \\
\midrule
NFR1 & Offline-First Operation & All core workflows (FR1--FR5) must function without connectivity; data persists locally and synchronizes when online \\
NFR2 & Cross-Platform Access & Single codebase accessible on desktop, tablet, and mobile via web browser; installable as PWA on mobile devices \\
NFR3 & Response Time & Page transitions under 2 seconds; form submissions under 1 second in offline mode \\
NFR4 & Sync Latency & Queued data synchronized within 30 seconds of connectivity restoration \\
NFR5 & Security & Role-based access control; Firebase Authentication; credential hashing for offline cache \\
NFR6 & Scalability & Architecture supports multiple industrial parks under a single deployment \\
NFR7 & Accessibility & WCAG 2.1 Level AA compliance for core interfaces \\
\bottomrule
\end{tabular}
\end{table}

\subsection{MoSCoW Prioritization}

The MoSCoW framework guided requirement prioritization and iteration planning. Offline availability served as the primary deciding factor. Table~\ref{tab:moscow} shows the results.

\begin{table}[H]
\centering
\caption{MoSCoW Feature Prioritization}
\label{tab:moscow}
\small
\begin{tabular}{p{2.5cm}p{5cm}p{4cm}}
\toprule
\textbf{Priority} & \textbf{Features} & \textbf{Offline Requirement} \\
\midrule
Must Have & Authentication, service requests, token system, announcements, complaint management & Full offline capability \\
Should Have & Asset management, notifications, billing, document management & Graceful degradation \\
Could Have & AI services, interactive map, OSS services, PDF export & Online-preferred \\
Won't Have (this phase) & External payment gateway, IoT integration, Amharic localization & Identified as future work \\
\bottomrule
\end{tabular}
\end{table}

\section{China-to-Ethiopia Technology Adaptation Framework}

\subsection{Chinese Smart Park Feature Analysis}

An examination of WeChat Work and DingTalk platforms revealed several feature categories common to Chinese smart parks: enterprise messaging and group communication, workflow approval routing, service request submission and tracking, payment processing (through WeChat Pay or Alipay), document management, attendance and access control, video conferencing, and analytics dashboards.

\subsection{Adaptation Dimensions}

Four dimensions structure the adaptation framework. The technological dimension carries the most weight --- the connectivity gap between Chinese and Ethiopian contexts is the primary constraint that reshapes everything else --- but organizational fit, cultural context, and economic realities each produce their own distinct adaptation requirements.

\textbf{Technological Dimension:} Connectivity demands the most fundamental adaptation. Chinese platforms assume persistent high-speed internet; the Ethiopian version replaces this with an offline-first architecture. Dependencies on Chinese ecosystem platforms (WeChat, Alipay, Alibaba Cloud) are replaced by globally accessible alternatives (Firebase, web standards).

\textbf{Organizational Dimension:} IPDC's organizational structure and service processes do not mirror those of Chinese park management. The adaptation translates Chinese workflow patterns into IPDC's three-role structure (tenant, admin, operations) and aligns approval workflows with the institution's existing hierarchies.

\textbf{Cultural Dimension:} Interface design was adjusted to fit local expectations. An English-language interface reflects the international business environment of Ethiopian industrial parks, while Amharic localization has been identified as a future enhancement. The digital service token concept is an adaptation of Chinese integrated payment, designed for a context where external payment gateway partnerships have not yet been established.

\textbf{Economic Dimension:} Cost constraints shaped technology choices at every stage: PWA over native apps (eliminating app store fees and multi-platform development costs), Firebase's free tier over paid cloud services, and open-source tools across the full stack.

\subsection{Feature Mapping and Modification Matrix}

Table~\ref{tab:adaptation-matrix} maps each Chinese smart park feature to its adapted counterpart in the IPDC platform.

\begin{table}[H]
\centering
\caption{China-to-Ethiopia Feature Adaptation Matrix}
\label{tab:adaptation-matrix}
\small
\begin{tabular}{p{3.5cm}p{3.5cm}p{5cm}}
\toprule
\textbf{Chinese Feature} & \textbf{IPDC Adaptation} & \textbf{Modification Rationale} \\
\midrule
WeChat Work messaging & Announcement system & Structured notifications replace real-time chat; works offline \\
DingTalk approval workflows & Service request pipeline & Adapted to tenant$\rightarrow$admin$\rightarrow$operations flow \\
WeChat Pay / Alipay & Digital service tokens & Internal system avoids external API dependency \\
Native mobile apps & Progressive Web App & Single codebase, no app store, offline-capable \\
Cloud-only architecture & Offline-first with Firebase & Dual-layer local storage + cloud sync \\
Real-time IoT monitoring & Asset management module & Manual tracking as IoT infrastructure unavailable \\
Chinese language UI & English interface & Primary business language; Amharic as future work \\
Alibaba Cloud backend & Firebase (Google Cloud) & Global availability, built-in offline persistence \\
\bottomrule
\end{tabular}
\end{table}

\subsection{Offline-First Adaptation Strategy}

The offline-first adaptation represents the most architecturally consequential departure from the Chinese reference models. It is worth being clear that this is not simply a matter of enabling a ``save offline'' toggle; it requires rethinking how data flows through the entire system. The adaptation is built on three guiding principles:

\begin{enumerate}
\item \textbf{Local-first data storage:} Every piece of data the user reads comes from local storage (the Firestore cache or IndexedDB), while network requests happen in the background to keep things synchronized.
\item \textbf{Queue-based write operations:} When users create or modify data while offline, their operations are placed in an IndexedDB queue and executed against Firestore once connectivity is restored.
\item \textbf{Transparent status communication:} The interface continuously shows connectivity status, synchronization state, and any pending operations, so users always know exactly where their data stands.
\end{enumerate}

\section{System Architecture}

\subsection{High-Level Architecture Overview}

The platform is organized as a three-tier architecture adapted for offline-first operation. Figure~\ref{fig:architecture} provides an overview of the high-level structure.

\begin{figure}[H]
\centering
\resizebox{0.95\textwidth}{!}{%
\begin{tikzpicture}[
  tier/.style={rectangle, draw, fill=blue!5, minimum width=14cm, minimum height=2.4cm, rounded corners=4pt},
  component/.style={rectangle, draw, fill=white, text width=2.8cm, text centered, minimum height=0.9cm, font=\footnotesize, rounded corners=2pt},
  arrow/.style={-{Stealth[length=2.5mm]}, thick},
  label/.style={font=\footnotesize\bfseries, text=blue!70}
]

% Presentation Tier
\node[tier, fill=green!5] (pres) at (0,4.5) {};
\node[label] at (-6,5.4) {Presentation Tier};
\node[component] at (-4.2,4.5) {React Components\\(41 components)};
\node[component] at (0,4.5) {MUI Interface\\(Responsive)};
\node[component] at (4.2,4.5) {Service Worker\\(Workbox)};

% Service Tier
\node[tier, fill=yellow!5] (serv) at (0,1.5) {};
\node[label] at (-6,2.4) {Service Tier};
\node[component] at (-4.2,1.5) {Firebase Auth\\+ Offline Auth};
\node[component] at (0,1.5) {Firestore SDK\\+ Offline Cache};
\node[component] at (4.2,1.5) {FastAPI\\(AI Models)};

% Data Tier
\node[tier, fill=red!5] (data) at (0,-1.5) {};
\node[label] at (-6,-0.6) {Data Tier};
\node[component] at (-4.2,-1.5) {Cloud Firestore\\(12 collections)};
\node[component] at (0,-1.5) {Firebase Storage\\(Documents)};
\node[component] at (4.2,-1.5) {IndexedDB / Dexie\\(Local queue)};

\draw[arrow] (pres) -- (serv);
\draw[arrow] (serv) -- (data);

\end{tikzpicture}}
\caption{High-Level System Architecture}
\label{fig:architecture}
\end{figure}

The \textbf{Presentation Tier} comprises 41 React components across six modules (admin, common, layout, map, operations, tenant) and 15 page-level components. Material-UI supplies the responsive design system; Workbox handles service worker management for asset caching and offline navigation.

The \textbf{Service Tier} mediates between the user interface and persistent storage. Firebase Authentication handles online identity verification, backed by an offline authentication module that checks cached credentials stored in IndexedDB. The Firestore SDK provides transparent offline caching, and a custom sync service manages the explicit operation queue. Separately, the FastAPI backend serves three AI models through REST endpoints.

The \textbf{Data Tier} blends cloud and local storage. Cloud Firestore holds twelve collections constituting the complete data model. Firebase Storage retains uploaded documents and images. IndexedDB (through Dexie.js) maintains a local queue of pending operations and cached service request data for offline access.

\subsection{Offline-First Architecture Design}

The offline-first architecture relies on a dual-layer storage strategy. Figure~\ref{fig:offline-flow} shows data flow during a service request submission under both online and offline conditions.

\begin{figure}[H]
\centering
\resizebox{0.85\textwidth}{!}{%
\begin{tikzpicture}[
  box/.style={rectangle, draw, fill=blue!8, text width=2.8cm, text centered, minimum height=1cm, rounded corners=3pt, font=\footnotesize},
  decision/.style={diamond, draw, fill=yellow!15, text width=1.5cm, text centered, aspect=2, font=\footnotesize},
  arrow/.style={-{Stealth[length=2.5mm]}, thick},
]

\node[box] (user) at (0,0) {User submits service request};
\node[decision] (check) at (0,-2.5) {Online?};
\node[box, fill=green!10] (direct) at (5,-2.5) {Write directly to Firestore};
\node[box, fill=orange!10] (queue) at (-5,-2.5) {Save to IndexedDB\\+ Add to syncQueue};
\node[box] (cache) at (-5,-5) {Firestore local cache updated};
\node[decision] (reconnect) at (-5,-7.5) {Connectivity restored?};
\node[box, fill=green!10] (sync) at (-5,-10) {SyncService processes queue $\rightarrow$ Firestore};

\draw[arrow] (user) -- (check);
\draw[arrow] (check) -- node[above]{\scriptsize Yes} (direct);
\draw[arrow] (check) -- node[above]{\scriptsize No} (queue);
\draw[arrow] (queue) -- (cache);
\draw[arrow] (cache) -- (reconnect);
\draw[arrow] (reconnect) -- node[left]{\scriptsize Yes} (sync);
\draw[arrow] (reconnect.east) -- ++(2,0) |- node[near start, right]{\scriptsize No (wait)} (cache.east);

\end{tikzpicture}}
\caption{Offline-First Data Flow for Service Request Submission}
\label{fig:offline-flow}
\end{figure}

\textbf{Layer 1 --- Firebase Firestore Persistent Cache:} With \texttt{persistentLocalCache()} and \texttt{persistentMultipleTabManager()} enabled, Firestore maintains a local copy of every document the user has previously accessed. When offline, read operations draw from this cache, giving users seamless access to existing records without any network connection.

\textbf{Layer 2 --- Dexie.js IndexedDB Queue:} Write operations initiated while offline are stored in two IndexedDB tables managed by Dexie.js. The \texttt{serviceRequests} table holds complete request data locally; the \texttt{syncQueue} table records metadata for each pending operation (type, collection, document ID, timestamp). Once the \texttt{useOnlineStatus} hook detects restored connectivity, the sync service processes queued items and writes them to Firestore.

\subsection{PWA Architecture}

The PWA setup, managed through vite-plugin-pwa, employs three caching strategies:

\begin{itemize}
\item \textbf{Precaching:} JavaScript bundles, CSS stylesheets, HTML pages, icons, and image assets are cached when the service worker is first installed, ensuring the application shell loads even without network access.
\item \textbf{CacheFirst} (Google Fonts): Font resources are served from cache with a one-year expiration, avoiding unnecessary repeat downloads.
\item \textbf{NetworkFirst} (Firebase Storage, API endpoints): For dynamic content, the application tries the network first and falls back to cached versions during outages. Firebase Storage responses are cached for seven days.
\end{itemize}

The web app manifest declares the application as \texttt{"display": "standalone"}, enabling full-screen installation on mobile devices. It also specifies the application name (``IPDC-OSS Platform''), a theme color (\texttt{\#2563eb}), and icon assets at 192px and 512px resolutions.

\subsection{Authentication Architecture}

Authentication follows a dual-mode design so that users can still log in during connectivity disruptions. Figure~\ref{fig:auth-flow} illustrates the authentication flow.

\begin{figure}[H]
\centering
\resizebox{0.95\textwidth}{!}{%
\begin{tikzpicture}[
  box/.style={rectangle, draw, fill=blue!8, text width=3cm, text centered, minimum height=0.9cm, rounded corners=3pt, font=\footnotesize},
  decision/.style={diamond, draw, fill=yellow!15, text width=1.5cm, text centered, aspect=2, font=\footnotesize},
  arrow/.style={-{Stealth[length=2.5mm]}, thick},
]

\node[box] (login) at (0,0) {User enters credentials};
\node[decision] (online) at (0,-2) {Online?};
\node[box, fill=green!10] (firebase) at (5,-2) {Firebase Auth\\signIn()};
\node[box] (cache) at (5,-4) {Cache credentials to IndexedDB};
\node[box, fill=orange!10] (offline) at (-5,-2) {Verify against IndexedDB cache};
\node[decision] (valid) at (-5,-4) {Valid?};
\node[box, fill=green!10] (grant) at (0,-6) {Grant access\\(role-based routing)};
\node[box, fill=red!10] (deny) at (-8.5,-4) {Deny access};

\draw[arrow] (login) -- (online);
\draw[arrow] (online) -- node[above]{\scriptsize Yes} (firebase);
\draw[arrow] (firebase) -- (cache);
\draw[arrow] (cache) |- (grant);
\draw[arrow] (online) -- node[above]{\scriptsize No} (offline);
\draw[arrow] (offline) -- (valid);
\draw[arrow] (valid) |- node[near start, right]{\scriptsize Yes} (grant);
\draw[arrow] (valid) -- node[above]{\scriptsize No} (deny);

\end{tikzpicture}}
\caption{Dual-Mode Authentication Flow}
\label{fig:auth-flow}
\end{figure}

When online, authentication proceeds through Firebase Auth via \texttt{signInWithEmailAndPassword()}. After a successful login, credentials are hashed and stored in IndexedDB through the \texttt{cacheCredentials()} utility. When offline, the system retrieves cached credentials and checks the user's input against the stored hash. The \texttt{AuthContext} component orchestrates this dual-mode logic, switching automatically to offline authentication whenever it encounters network errors.

Three user roles --- \texttt{tenant}, \texttt{admin}, and \texttt{operations} --- determine which dashboard a user sees and which features they can access, managed through role-based conditional rendering in the React component tree.

\subsection{AI Integration Architecture}

Three machine learning models are served through a FastAPI backend as REST endpoints at \texttt{/api/v1/}:

\begin{enumerate}
\item \textbf{Service Classifier} (\texttt{/classify-service}): Accepts a request title and description; returns a predicted service type, priority level, estimated processing time in days, confidence score, and contextual recommendations.
\item \textbf{Predictive Maintenance} (\texttt{/predict-maintenance}): Takes asset attributes --- age, condition score, maintenance history --- and outputs a failure probability, risk level, estimated days until failure, and a recommended action. The output is designed to give operations teams something actionable rather than just a probability number.
\item \textbf{Park Recommendation} (\texttt{/recommend-parks}): Accepts a tenant profile (industry sector, employee count, investment capital, infrastructure needs); returns ranked park recommendations with match scores, pros and cons, and cost estimates. Scoring weights are: industry alignment 30\%, capacity 25\%, infrastructure 20\%, cost 15\%, location 10\%.
\end{enumerate}

AI features are designed with graceful degradation: when the FastAPI backend is unreachable during connectivity outages, the platform reverts to manual classification and disables recommendation features without blocking core workflows.

\section{Database Design}

\subsection{Firestore Collections Schema}

The Firestore database is organized into twelve collections. Table~\ref{tab:firestore-schema} documents the primary collections with their key fields and purposes.

\begin{table}[H]
\centering
\caption{Firestore Collections Schema}
\label{tab:firestore-schema}
\small
\begin{tabular}{p{3cm}p{5cm}p{4cm}}
\toprule
\textbf{Collection} & \textbf{Key Fields} & \textbf{Purpose} \\
\midrule
users & uid, email, displayName, role, parkId, createdAt & User profiles and roles \\
serviceRequests & id, tenantId, title, description, serviceType, priority, status, assignedTo & Core service workflow \\
ossServiceRequests & id, tenantId, serviceCategory, documents[], status & OSS-specific requests \\
announcements & id, title, content, targetAudience, priority, expiresAt, isActive & System announcements \\
complaints & id, requestId, tenantId, category, evidence[], status, resolution & Complaint tracking \\
tokenAccounts & tenantId, balance, tier, totalEarned, totalSpent & Token balances \\
tokenTransactions & id, tenantId, type, amount, description, timestamp & Transaction ledger \\
invoices & id, tenantId, lineItems[], total, status, pdfUrl & Billing records \\
paymentTransactions & id, tenantId, amount, method, referenceNumber & Payment tracking \\
assets & id, name, category, condition, purchaseCost, currentValue & Park asset inventory \\
maintenanceRecords & id, assetId, type, scheduledDate, completedDate, cost & Maintenance history \\
notifications & id, userId, title, message, read, createdAt & User notifications \\
\bottomrule
\end{tabular}
\end{table}

\subsection{IndexedDB Schema (Dexie.js)}

The local IndexedDB database, named \texttt{IPDCDatabase}, holds two tables tailored for offline operations:

\begin{itemize}
\item \texttt{serviceRequests}: Stores complete request documents for offline access. Indexed on: \texttt{id}, \texttt{tenantId}, \texttt{status}, \texttt{priority}, \texttt{serviceType}, \texttt{createdAt}.
\item \texttt{syncQueue}: Stores pending write operations with auto-incrementing IDs. Indexed on: \texttt{++id}, \texttt{type}, \texttt{collection}, \texttt{status}, \texttt{createdAt}.
\end{itemize}

The sync queue recognizes four status values: \texttt{pending} (awaiting sync), \texttt{syncing} (currently being processed), \texttt{synced} (successfully written to Firestore), and \texttt{failed} (sync attempted but unsuccessful; will be retried).

\section{User Interface Design}

\subsection{Design Principles}

The interface follows Material Design guidelines through the MUI component library, enhanced with a custom theme. A primary color (\texttt{\#2563eb}, professional blue) and secondary color (\texttt{\#059669}, confident green) establish a visual identity suited to an institutional platform. Status-specific colors give instant feedback: amber for pending, blue for approved, purple for in-progress, green for completed, and red for rejected.

Responsive breakpoints at 600px and 960px ensure the platform works across device categories. The \texttt{useDeviceDetection} hook dynamically adjusts layout density and navigation patterns based on the device in use.

\subsection{Navigation and Role-Based Views}

After authentication, the routing system directs each user to a role-appropriate dashboard:

\begin{itemize}
\item \textbf{Tenant Dashboard:} Service request submission, request tracking, token management, complaint filing, park recommendations.
\item \textbf{Admin Dashboard:} Request management (across all tenants), announcement creation, tenant account management, complaint resolution, billing oversight, asset management.
\item \textbf{Operations Dashboard:} Service fulfillment, maintenance scheduling, asset inspection, performance metrics.
\end{itemize}

\subsection{Status Indicators}

The \texttt{StatusBar} component keeps users informed about system state at all times, showing:
\begin{itemize}
\item Online/offline connectivity status with a color-coded indicator;
\item Device type (desktop, tablet, or mobile);
\item Number of pending sync operations;
\item Timestamp of the last successful synchronization.
\end{itemize}

This transparency is essential for an offline-first system, where users need to know whether their actions have been synchronized to the server or are still waiting in a local queue.

\section{Chapter Summary}

The analysis and design work described in this chapter produced several concrete outputs: twelve functional and seven non-functional requirement groups grounded in stakeholder input, a four-dimension adaptation framework translating Chinese platform patterns to the Ethiopian context, and a three-tier architecture built from 41 React components, Firebase and FastAPI services, and a dual-layer local storage strategy. Of these, the dual-mode authentication and queue-based synchronization mechanisms arguably matter most to end users --- they are what make the platform feel continuous and trustworthy rather than intermittent and fragile. It is also worth noting that many of the design decisions recorded here were not obvious at the outset; they emerged iteratively as the constraints of the Ethiopian operating environment became clearer. How those decisions translate into working software, and what the evaluation reveals about them, is the subject of the next chapter.
